\section{Introduction}
\label{sec:intro}

In clusterwise linear regression each observation follows one of $K$ linear models, and the
model that generated it is not observed. For units $(x_i,y_i)$, $i=1,\dots,n$, with $x_i\in\R^{d}$
including an intercept,
\begin{equation}
y_i=x_i^{\top}\theta^{0}_{z_i}+\sigma_{z_i}e_i,\qquad z_i\in\{1,\dots,K\},
\label{eq:model}
\end{equation}
where $z_i$ is the unobserved label of unit $i$, $\theta^{0}_k\in\R^{d}$ and $\sigma_k>0$ are
the coefficients and noise scale of group $k$, and the $e_i$ are independent standard errors.
The model goes back to \citet{quandt1972switching} and \citet{spath1979algorithm}. It is
usually fitted by least squares over partitions or by maximum likelihood for a mixture of
regressions \citep{desarbo1988maximum,deveaux1989mixtures}.

Both fits break down when a fraction of the data does not follow any of the $K$ models: a few
gross outliers in the response can pull a fitted line far from every group, or be fitted as a
group of their own. The standard remedy is trimming. TCLUST-REG
\citep{garciaescudero2010robust} discards the fraction $\alpha$ of units that fit worst and
fits the model to the rest; the trimmed likelihood estimator (TLE) of \citet{neykov2007robust} and the
trimmed cluster-weighted model \citep{garciaescudero2017robust} do the same within other
likelihoods. Each needs the trimming level $\alpha$, which the analyst rarely knows. If $\alpha$
is smaller than the fraction of outliers, the fit breaks down; if it is much larger than that
fraction, genuine units are discarded, and a small group can be trimmed away altogether.

Several ways around this choice have been proposed. One trims at a deliberately high level
and then reinstates the units that are well fitted, a reweighting step proposed for robust
clustering by \citet{dotto2018reweighting}; it needs a level that is high enough. Another
monitors the fit over a grid of levels and leaves the choice to the analyst
\citep{torti2019assessing,riani2008fitting}, or scans the number of groups and the level
jointly by trimmed likelihood curves \citep{garciaescudero2011exploring,cerioli2018finding}.
A third adapts the trimming within each group \citep{chang2020componentwise}, and a fourth
replaces trimming by bounded loss functions \citep{bai2012robust}. Mixtures with a mean-shift
parameter for each unit \citep{yu2017new}, with a uniform noise component
\citep{banfield1993model}, or with contaminated normal errors
\citep{punzo2017robust,mazza2020mixtures} model the outliers instead of trimming them. Multiple-model
fitting from random minimal samples, from RANSAC
\citep{fischler1981random} to J-linkage \citep{toldo2008robust}, T-linkage
\citep{magri2014tlinkage} and preference analysis \citep{magri2015robust}, treats a related
problem with the noise scale as a user-set threshold, and does not ask whether what it
discards is a missed model.

This paper makes two proposals. The first concerns what a robust fit throws away. When one
group is small, a trimming method at a generous level can trim it away with the outliers, and a
fit that flags units by their residuals can spend two lines on the large group and flag the small
one; our simulations show both happening to a group of $20\%$ of the data even without outliers.
The \emph{group-recovery step} (Algorithm~\ref{alg:recovery}) searches the trimmed or flagged
units for a line around which many of them are concentrated, and keeps the line if the
likelihood of Gaussian groups plus uniform noise prefers the fit that contains it. The step
needs only a fit and its flagged set, so it can follow TCLUST-REG with reweighting as well as
the level-free procedure introduced next; it is what makes a generous level safe with
unequal groups. Trimmed likelihood curves ask the same question by a joint scan; the step
asks it locally and post hoc (Section~\ref{sec:recovery}).

The second proposal is a flagging procedure that needs no level, the exact-subsample flagging
procedure, ESF (Algorithm~\ref{alg:adaptive}). A small random subsample that contains no
outlier, fitted by least squares over all its partitions, gives a clean fit of the model, and
that fit identifies the outliers of the whole sample as the units far from every fitted line.
ESF therefore draws many small subsamples, fits each one exactly, scores each fit by a bounded
loss on the units not yet flagged, and flags the units that lie far from the best fit; later
subsamples are drawn only from units not yet flagged. The fit ends with the reweighting step of
robust regression \citep{rousseeuw1987robust}, and outlying covariates are flagged by a
separate screen. A unit is flagged when its residual exceeds a fixed multiple of a robust
scale; this cut-off takes the place of the trimming level, and it does not depend on how many
outliers there are. Two constants stand in for the level: the subsample size $m$, set from a
lower bound on the smallest group proportion, and the cap of $n/3$ on the flagged set beyond
which the recovery step is not attempted; Section~\ref{sec:constants} gives a default rule for
each. Subsampling to find a clean subset is the idea behind least median of squares and least
trimmed squares \citep{rousseeuw1984least,rousseeuw2006computing}, behind robust linear
grouping \citep{garciaescudero2009robust} and behind the random starts of TCLUST-REG. Here
each subsample contains several units of every group and is solved to global optimality, so
a clean subsample yields the clusterwise least-squares fit of that subsample and not a fit to
$d$ points. The exact fit is the problem that the operations research literature solves on
whole samples \citep{bertsimas2007classification,carbonneau2011globally,carbonneau2012extensions,carbonneau2014globally,fois2026clusterwise,bagirov2013nonsmooth,joki2020clusterwise,park2017algorithms},
without treating outliers. Here it is applied about a hundred times per data set to subsamples
of a dozen units, where a plain depth-first branch and bound is one to two orders of magnitude
faster than a mixed-integer programme. The accuracy of ESF does not rest on the exactness:
with ten starts of hard alternation in place of the exact fit of each subsample, ESF alone was
within $0.01$ of the exact fit in most cells, ahead by $0.025$ in one cell and behind by
$0.02$ and $0.04$ with four groups (Supplement, Table~S14). The exact fit removes instead the
dependence of a subsample's fit on its starting point.

The place of ESF among the trimming methods should be stated at once. In the simulation
study, TCLUST-REG at the fixed level $0.30$ with reweighting and the recovery step matched the
mean accuracy of ESF with up to $20\%$ of outliers and had no cell below $0.85$, where ESF had
two; the level $0.25$ with equal group weights did better still; and the levels $0.35$ and
$0.40$ beat ESF beyond $25\%$ of outliers (Table~\ref{tab:level}). The level $0.35$ came
closest to serving both ranges; up to $20\%$ of outliers it fell behind ESF only with three and
four groups, and at $35\%$ it fell behind in four designs with two groups. ESF is
therefore not the more accurate method; its case is that it needs no level and holds up with
three and four groups, where every fixed level with estimated group weights loses accuracy.

The contributions are the following.
\begin{enumerate}[label=(\roman*),leftmargin=2.2em,itemsep=2pt]
\item The group-recovery step, which can follow any trimming or flagging method
(Section~\ref{sec:recovery}).
\item ESF (Section~\ref{sec:method}), with a comparison of exact solvers for the small
clusterwise least-squares problems and an ablation of the exact fit (Section~\ref{sec:solvers}).
\item Three partial results (Section~\ref{sec:theory}): the exact probability that a stage of
the sequential draw contains a clean subsample, the large-sample limit of the final reweighting
step and of the fraction it flags, and the value of the two operations of the recovery step in
the population.
\item A simulation study on $7{,}750$ data sets (Section~\ref{sec:simulation}).
After TCLUST-REG at a generous level with reweighting, the recovery step repaired nearly all
failures with unequal or small groups but not those with three or four groups. ESF followed by
the recovery step was within $0.03$ of the better of the levels $0.25$ and $0.40$, chosen after
the run, in $83$ of $89$ cells with up to a fifth of outliers, and was the less robust choice
beyond a quarter.
\item Three data sets (Section~\ref{sec:realdata}): the fishery and tone perception data, and taxi
fares from an airport, where the tariff of every trip is recorded and so the true group of every
unit is known.
\end{enumerate}
No method in the study is best everywhere; Section~\ref{sec:discussion} says which to use
when. The constants of ESF and of the recovery step were the same in every experiment.
